\documentclass[aspectratio=169,11pt]{beamer}
\usetheme{metropolis}
\usepackage{fontspec}
\usepackage{polyglossia}
\setmainlanguage{ukrainian}
\setmainfont{DejaVuSans}[Extension=.ttf,UprightFont=*,BoldFont=*-Bold,ItalicFont=*-Oblique,BoldItalicFont=*-BoldOblique]
\setsansfont{DejaVuSans}[Extension=.ttf,UprightFont=*,BoldFont=*-Bold,ItalicFont=*-Oblique,BoldItalicFont=*-BoldOblique]
\usepackage{booktabs,tabularx,array,colortbl,xcolor}
\definecolor{sgreen}{HTML}{E2F0D9}
\definecolor{wred}{HTML}{FBE5D6}
\definecolor{oblue}{HTML}{DEEBF7}
\definecolor{tyellow}{HTML}{FFF2CC}
\newcolumntype{Y}{>{\raggedright\arraybackslash}X}
\setbeamerfont{frametitle}{size=\large}
\setbeamertemplate{itemize items}{\textbullet}
\setlength{\leftmargini}{1.2em}

\newcommand{\swotslide}[5]{%
\begin{frame}{#1 --- SWOT}
\footnotesize
\renewcommand{\arraystretch}{1.05}
\begin{tabularx}{\textwidth}{|Y|Y|}
\hline
\cellcolor{sgreen}\textbf{S — сильні сторони} & \cellcolor{wred}\textbf{W — слабкі сторони}\\
\hline
\cellcolor{sgreen!40}\begin{itemize}#2\end{itemize} &
\cellcolor{wred!40}\begin{itemize}#3\end{itemize}\\
\hline
\cellcolor{oblue}\textbf{O — можливості} & \cellcolor{tyellow}\textbf{T — загрози}\\
\hline
\cellcolor{oblue!40}\begin{itemize}#4\end{itemize} &
\cellcolor{tyellow!40}\begin{itemize}#5\end{itemize}\\
\hline
\end{tabularx}
\end{frame}}

\title{Вибір та обґрунтування теми командного проєкту}
\subtitle{Інформаційні технології, Модуль 2 — Етап 1}
\author{Петров Богдан \and Мекшаков Артемій}
\institute{Група МІ-41}
\date{21 жовтня 2026}

\begin{document}
\maketitle

\begin{frame}{Команда та обмеження}
\begin{columns}[T]
\column{0.48\textwidth}
\textbf{Команда}
\begin{itemize}
  \item Петров Богдан — Team Lead, Backend (NestJS), БД
  \item Мекшаков Артемій — Product Owner, Frontend (Next.js), QA
\end{itemize}
\column{0.48\textwidth}
\textbf{Критерії відбору тем}
\begin{itemize}
  \item реалістичний обсяг для 2 осіб до 18.11
  \item стек: TypeScript, Next.js, NestJS, PostgreSQL
  \item реальна потреба, доступна аудиторія
  \item придатність до тестування та API-документування
\end{itemize}
\end{columns}
\end{frame}

\begin{frame}{Три варіанти теми}
\begin{enumerate}
  \item \textbf{StudyFlow} — планувальник навчальних дедлайнів для студентських груп
  \item \textbf{SplitMate} — облік спільних витрат і розрахунок боргів
  \item \textbf{QuickPoll} — миттєві опитування та голосування для занять
\end{enumerate}
\end{frame}

\begin{frame}{Тема 1. StudyFlow}
\small
\begin{tabularx}{\textwidth}{>{\bfseries}p{3cm} Y}
Що & Спільний календар дедлайнів, лабораторних і командних задач групи\\[3pt]
Для кого & Студенти, навчальні групи/підгрупи, старости\\[3pt]
Проблема & Дедлайни розпорошені між Classroom, чатами й усними домовленостями; немає єдиної картини тижня\\[3pt]
Можливості & Групи та ролі; предмети й дедлайни; канбан задач; календар; e-mail-нагадування; прогрес групи\\
\end{tabularx}
\end{frame}

\swotslide{StudyFlow}{
  \item команда = цільова аудиторія
  \item простий CRUD + ролі + календар, ідеально для Next/Nest
  \item чітка реляційна модель, легко тестувати
}{
  \item низька технологічна новизна
  \item цінність залежить від ручного внесення даних
  \item нагадування потребують фонових задач
}{
  \item інтеграція з Google Calendar / Classroom
  \item пілот у власній групі МІ-41
  \item масштабування на факультет
}{
  \item конкуренція з Notion, Google Calendar, Trello
  \item ризик низького залучення
  \item залежність від e-mail-сервісу
}

\begin{frame}{Тема 2. SplitMate}
\small
\begin{tabularx}{\textwidth}{>{\bfseries}p{3cm} Y}
Що & Облік спільних витрат і автоматичний розрахунок «хто кому винен»\\[3pt]
Для кого & Студенти в гуртожитку, компанії в поїздках, співмешканці\\[3pt]
Проблема & Розрахунки губляться в чатах, підрахунок балансу вручну трудомісткий і спричиняє конфлікти\\[3pt]
Можливості & Групи; витрати з розподілом; баланси; мінімізація переказів; історія, експорт CSV\\
\end{tabularx}
\end{frame}

\swotslide{SplitMate}{
  \item формалізована бізнес-логіка, зручна для unit-тестів
  \item мало сутностей, швидкий MVP
  \item реальна побутова потреба
}{
  \item робота з грошима: точність, округлення
  \item обмежений простір для UI
  \item без мобільного застосунку менш зручний
}{
  \item валюти, розпізнавання чеків, банківські API
  \item преміум-функції
}{
  \item зрілі конкуренти: Splitwise, Tricount
  \item низька лояльність користувачів
  \item довіра до фінансових даних
}

\begin{frame}{Тема 3. QuickPoll}
\small
\begin{tabularx}{\textwidth}{>{\bfseries}p{3cm} Y}
Що & Швидкі опитування й голосування за кодом/QR з результатами в реальному часі\\[3pt]
Для кого & Викладачі та студенти на заняттях, організатори заходів\\[3pt]
Проблема & Потрібно швидко зібрати думку аудиторії без реєстрації та складних інструментів\\[3pt]
Можливості & Типи питань; доступ за кодом; real-time результати (WebSocket); діаграми; експорт\\
\end{tabularx}
\end{frame}

\swotslide{QuickPoll}{
  \item ефектне live-демо
  \item без реєстрації для учасників
  \item корисний самим викладачам
}{
  \item WebSocket ускладнює архітектуру і тести
  \item стабільність при багатьох підключеннях
  \item анонімність — накрутка голосів
}{
  \item інтеграція з LMS, конференції
  \item гейміфіковані вікторини
}{
  \item Mentimeter, Kahoot, Slido, Google Forms
  \item ризик не стабілізувати real-time до дедлайну
  \item низька потреба: є звичні інструменти
}

\begin{frame}{Зважене порівняння}
\small
\renewcommand{\arraystretch}{1.2}
\begin{tabularx}{\textwidth}{Y c c c c}
\toprule
\textbf{Критерій} & \textbf{Вага} & \textbf{StudyFlow} & \textbf{SplitMate} & \textbf{QuickPoll}\\
\midrule
Реалізованість 2 особами до 18.11 & 0,30 & 5 & 5 & 3\\
Відповідність вимогам курсу & 0,25 & 5 & 4 & 3\\
Доступ до аудиторії & 0,20 & 5 & 4 & 3\\
Потенціал розвитку & 0,15 & 4 & 3 & 4\\
Конкурентне середовище (5 — слабке) & 0,10 & 3 & 2 & 1\\
\midrule
\textbf{Зважена оцінка} & & \textbf{4,65} & \textbf{3,95} & \textbf{2,95}\\
\bottomrule
\end{tabularx}
\end{frame}

\begin{frame}{Обраний варіант: StudyFlow}
\begin{itemize}
  \item \textbf{Реалізованість:} відомі патерни (auth, ролі, CRUD, календар), без ризикового real-time
  \item \textbf{Вимоги курсу:} змістовна схема БД та UML, 20+ unit-тестів на бізнес-логіку, Swagger
  \item \textbf{Аудиторія:} перевірка гіпотез на власній групі вже в семестрі
  \item \textbf{Рівне навантаження:} однаково вагомі backend і frontend частини
  \item \textbf{Загрози керовані:} вузька спеціалізація на навчальному процесі; ізольований модуль нагадувань
\end{itemize}
\vspace{4pt}
SplitMate — резервна альтернатива; QuickPoll — відхилено.
\end{frame}

\begin{frame}{Наступні кроки (етап 2, до 27.10)}
\begin{itemize}
  \item Lean Canvas для StudyFlow
  \item Карта продукту
  \item Дошка проєкту на платформі управління (Trello / Jira)
\end{itemize}
\vfill
\centering \large Дякуємо за увагу!
\end{frame}

\end{document}
